\chapter{State execution and implementation}
\source{Sample Q1/Q2/Q4, pp. 1/2/4; State and StateMonad lectures.}
Use this representation and the reference operations:
\par\noindent\begin{minipage}{\linewidth}
\begin{lstlisting}
newtype State s a = S {runState :: s -> (a,s)}
\end{lstlisting}
\end{minipage}\par
In every trace, keep a column for the returned value and another for
state after the step. Report \code{(result, finalState)} in that order.

\task{5.1. Evaluate and infer}
\begin{enumerate}[label=(\alph*)]
\item \code{runState (get >>= \n -> put (n*2) >> get) (5 :: Int)}
\answer{[Step table, result, type.]}
\item \code{runState (put False >> get) True}
\answer{[Step table, result, type.]}
\item \code{runState (modify reverse >> get) "abc"}
\answer{[Step table, result, type.]}
\item \code{runState ((,) <$> get) 'q'}
\answer{[Returned function and final state; then apply the returned function to True.]}
\item Using \code{bundle} from chapter 4, evaluate
the following with Boolean functions as state:
\par\noindent\begin{minipage}{\linewidth}
\begin{lstlisting}
runState (bundle [get, put not >> get]) id
\end{lstlisting}
\end{minipage}\par
\answer{[Returned function list and final function; apply each to False.]}
\end{enumerate}

\task{5.2. Construct State actions from types}
\code{swapStore} returns the reversed pair while leaving the stored pair
unchanged. \code{projectStore f x} returns \code{f currentState x}, also leaving
state unchanged. Use only the abstract State interface; do not construct \code{S}.
\par\noindent\begin{minipage}{\linewidth}
\begin{codeanswer}
swapStore :: State (a,b) (b,a)
swapStore = undefined  -- YOUR CODE

projectStore :: (a -> b -> c) -> b -> State a c
projectStore = undefined  -- YOUR CODE
\end{codeanswer}
\end{minipage}\par
\answer{[For each: stored-state type; returned-value type.]}

\task{5.3. Ordered updates}
\code{fresh} returns the old integer counter and increments it once.
\code{takeLabels n} collects \code{n} fresh labels; \code{n <= 0} leaves state
unchanged. \code{pop} returns and removes the head of list state, or returns
\code{Nothing} without changing empty state. No library traversal helpers.
\par\noindent\begin{minipage}{\linewidth}
\begin{codeanswer}
fresh :: State Int Int
fresh = undefined  -- YOUR CODE

takeLabels :: Int -> State Int [Int]
takeLabels = undefined  -- YOUR CODE

pop :: State [a] (Maybe a)
pop = undefined  -- YOUR CODE
\end{codeanswer}
\end{minipage}\par

\task{5.4. Label a tree}
Use \code{fresh} to attach consecutive labels to leaves, left subtree before
right, advancing the counter once per leaf. Do not construct \code{S} or call
\code{runState} inside your implementation.
\par\noindent\begin{minipage}{\linewidth}
\begin{codeanswer}
data Tree a = Leaf a | Branch (Tree a) (Tree a)
labelTree :: Tree a -> State Int (Tree (a,Int))
labelTree = undefined  -- YOUR CODE
\end{codeanswer}
\end{minipage}\par

\task{5.5. Function-valued state}
Return the currently stored Boolean function and replace the stored function by
\code{old . not}. Use the abstract State interface.
\par\noindent\begin{minipage}{\linewidth}
\begin{codeanswer}
snapshotFunction :: State (Bool -> Bool) (Bool -> Bool)
snapshotFunction = undefined  -- YOUR CODE
\end{codeanswer}
\end{minipage}\par
\answer{[How would you check both functions without Show or Eq instances for functions?]}

\task{5.6. Changed type question}
Infer \code{\f x -> f <$> (put x >> get)}. Describe the symbolic result of
running this action from a matching initial state \code{s0}.
\answer{[Most general type; result/state pair in terms of f, x, and s0.]}
